\begingroup
\small
\begin{thebibliography}{99}
\addcontentsline{toc}{section}{References}

\bibitem{Repo}
ProveIt Contributors,
\emph{Polylogarithms and their Arithmetic Bridges},
canonical manuscript and editorial records, revision
\texttt{\sourcecommit}, 11 October 2026.
\url{https://github.com/VladimirReshetnikov/ProveIt/tree/f7e35be357654198c1546fd40f363694c5bcc2cc/Analysis/Polylogarithms/docs/manuscript}.

\bibitem{Higher}
ProveIt research continuation,
\emph{Higher Digamma Powers and Ordered Gamma Sums}, 11 October 2026.
Incoming archive at the same revision; especially the all-power
Mittag--Leffler theorem, the finite harmonic-word algorithm, and the
differential-polynomial, common-shift, and higher-moment questions.
\url{https://github.com/VladimirReshetnikov/ProveIt/blob/f7e35be357654198c1546fd40f363694c5bcc2cc/docs/incoming/ProveIt_Higher_Digamma_Powers_2026-10-11.zip}.

\bibitem{Fifth}
ProveIt research continuation,
\emph{Canonical Digamma Polynomials and the Fifth Harmonic Bridge:
Laurent Coefficients, Polygamma Primitives, and Finite Reconstruction
of Even-Tail Moments}, 11 October 2026.
Incoming archive at the same revision; Sections 4, 7, and 9,
including Conjecture \texttt{conj:firstq}.
\url{https://github.com/VladimirReshetnikov/ProveIt/blob/f7e35be357654198c1546fd40f363694c5bcc2cc/docs/incoming/ProveIt_Canonical_Digamma_Fifth_Moment_2026-10-11.zip}.

\bibitem{Mixed}
ProveIt research continuation,
\emph{Mixed Resolvent Products, Higher Digamma Powers, and Stieltjes
Collisions: Exact Identities and Normalized Finite Parts},
11 October 2026. Incoming archive at the same revision.
\url{https://github.com/VladimirReshetnikov/ProveIt/blob/f7e35be357654198c1546fd40f363694c5bcc2cc/docs/incoming/ProveIt_Mixed_Resolvents_Harmonic_Powers_2026-10-11.zip}.

\bibitem{Continuous}
ProveIt research continuation,
\emph{Continuous-Depth Lerch Calculus: Gamma Kernels, All-Order Harmonic
Identities, Stieltjes Convolutions, and Anchored Primitives},
10 October 2026. Incoming archive at the same revision; the centered
reflection theorem, Gamma convolution, and final research question Q8.
\url{https://github.com/VladimirReshetnikov/ProveIt/blob/f7e35be357654198c1546fd40f363694c5bcc2cc/docs/incoming/ProveIt_Continuous_Depth_Lerch_2026-10-10.zip}.

\bibitem{Tilted}
ProveIt research continuation,
\emph{Tilted Lerch Convolutions and Binary Harmonic Cancellation:
Confluent Zeta Identities, Independent Spectral Jets, Gamma Primitives,
and Mixed-Parity Exact Evaluations}, 10 October 2026.
Incoming archive at the same revision.
\url{https://github.com/VladimirReshetnikov/ProveIt/blob/f7e35be357654198c1546fd40f363694c5bcc2cc/docs/incoming/ProveIt_Tilted_Lerch_Binary_Cancellation_2026-10-10.zip}.

\bibitem{Algebraic}
ProveIt research continuation,
\emph{Algebraic Gauss Residues and Centered Harmonic Reflection},
incoming archive at the same revision.
\url{https://github.com/VladimirReshetnikov/ProveIt/blob/f7e35be357654198c1546fd40f363694c5bcc2cc/docs/incoming/Algebraic_Gauss_Centered_Harmonic_Reflection.zip}.

\bibitem{DLMFGamma}
NIST Digital Library of Mathematical Functions,
Chapter 5, especially \S\S5.7 and 5.15,
Gamma expansions, polygamma recurrences, and special values.
\url{https://dlmf.nist.gov/5.7};
\url{https://dlmf.nist.gov/5.15}.

\bibitem{DLMFZeta}
NIST Digital Library of Mathematical Functions,
\S25.11, Hurwitz zeta function.
\url{https://dlmf.nist.gov/25.11}.

\bibitem{DLMFLerch}
NIST Digital Library of Mathematical Functions,
\S25.14, Lerch's transcendent.
\url{https://dlmf.nist.gov/25.14}.

\bibitem{DLMFEuler}
NIST Digital Library of Mathematical Functions,
\S3.9(ii), Euler's transformation of series.
\url{https://dlmf.nist.gov/3.9#ii}.

\bibitem{IKZ}
K. Ihara, M. Kaneko, and D. Zagier,
Derivation and double shuffle relations for multiple zeta values,
\emph{Compositio Mathematica} \textbf{142} (2006), 307--338.
\url{https://doi.org/10.1112/S0010437X0500182X}.
Author-hosted full text:
\url{https://www2.math.kyushu-u.ac.jp/~mkaneko/papers/34Derivation_double_shuffle.pdf}.

\bibitem{RamanujanFourier}
D. Chakrabarti and G. K. Srinivasan,
\emph{On a remarkable formula of Ramanujan}, 2012.
\url{https://arxiv.org/abs/1206.5122}.
The Gamma-product Fourier normalization used here is also derived
directly from the beta integral in Section~\ref{lj:sec:main}.

\end{thebibliography}
\endgroup
